\subsection*{۱}
ضرب داخلی پرس‌وجو و کلید را می‌نویسیم:
\begin{flalign*}
q\cdot k = \sum_{i=1}^{d_k} q_i k_i &&
\end{flalign*}
چون $q_i, k_i \sim \mathcal{N}(0,1)$ و مستقل‌اند، $E[q_i]=E[k_i]=0$ و $\mathrm{Var}(q_i)=\mathrm{Var}(k_i)=1$. از این رو:
\begin{flalign*}
E[q_i k_i] &= E[q_i]\,E[k_i] = 0 \\
\mathrm{Var}(q_i k_i) &= E[q_i^2 k_i^2] - \big(E[q_i k_i]\big)^2 = E[q_i^2]\,E[k_i^2] = 1 &&
\end{flalign*}
به دلیل استقلال جمله‌ها واریانس مجموع برابر مجموع واریانس‌هاست:
\begin{flalign*}
\mathrm{Var}(q\cdot k) = \sum_{i=1}^{d_k} \mathrm{Var}(q_i k_i) = d_k &&
\end{flalign*}

استدلال گرادیان: بدون ضریب $\frac{1}{\sqrt{d_k}}$ ورودی‌های \lr{softmax} واریانس $d_k$ و انحراف معیار $\sqrt{d_k}$ دارند؛ پس با بزرگ شدن $d_k$ لاجیت‌ها بزرگ می‌شوند و خروجی \lr{softmax} به برداری تقریباً \lr{one-hot} (اشباع) میل می‌کند. ژاکوبین \lr{softmax} برابر است با:
\begin{flalign*}
\frac{\partial p_i}{\partial z_j} = p_i(\delta_{ij} - p_j) &&
\end{flalign*}
وقتی یک $p_i \approx 1$ و بقیه $\approx 0$ باشند، $p_i(1-p_i)\approx 0$ و همهٔ مشتق‌ها به صفر میل می‌کنند؛ یعنی گرادیانِ عبوری از \lr{softmax} محو می‌شود. ضرب در $\frac{1}{\sqrt{d_k}}$ واریانس را مستقل از $d_k$ روی ۱ نگه می‌دارد و از این اشباع جلوگیری می‌کند.

\subsection*{۲}
نخست توجه کنیم هزینهٔ پارامتری دو حالت یکسان است: $h$ سر با بعد $d_k = d_{\text{model}}/h$ روی هم $h\cdot d_{\text{model}}\cdot \frac{d_{\text{model}}}{h} = d_{\text{model}}^2$ پارامتر در هر نگاشت دارند، دقیقاً مانند یک سر با $d_k = d_{\text{model}}$. پس برتری از تعداد پارامتر نمی‌آید.

یک سرِ تنها تنها یک توزیع توجه تولید می‌کند؛ یعنی خروجی هر پرس‌وجو صرفاً یک میانگین وزن‌دارِ (محدب) از بردارهای \lr{value} است. اگر یک توکن لازم باشد اطلاعات را از چند موقعیتِ مختلف و به دلایل متفاوت گرد آورد، این میانگین‌گیریِ یگانه آن‌ها را در هم می‌آمیزد و تمایزشان از میان می‌رود.

در مقابل، $h$ سر با نگاشت‌های مستقل $W_i^Q, W_i^K, W_i^V$ داده را به $h$ زیرفضای متفاوت می‌برند و $h$ توزیع توجهِ مستقل می‌سازند. بنابراین مدل می‌تواند هم‌زمان به موقعیت‌های مختلف و روابط گوناگون (نحوی، معنایی، مکانی و ...) توجه کند و سپس نتایج را الحاق نماید. همین توانایی نمایش هم‌زمانِ چند الگوی توجه در زیرفضاهای متفاوت بیانگری را بالا می‌برد و با یک سرِ بزرگ که ناگزیر به تولید تنها یک توزیع است حاصل نمی‌شود.

\subsection*{۳}
\subsubsection*{الف}
از تعریف کدگذاری، ضرب داخلی را باز می‌کنیم:
\begin{flalign*}
p_a^T p_b = \sum_{i=1}^{m}\big[\sin(\omega_i a)\sin(\omega_i b) + \cos(\omega_i a)\cos(\omega_i b)\big] &&
\end{flalign*}
با قرار دادن $x=\omega_i a$ و $y=\omega_i b$ در رابطهٔ $\cos(x-y)=\cos x\cos y + \sin x\sin y$:
\begin{flalign*}
p_a^T p_b = \sum_{i=1}^{m}\cos\big(\omega_i(a-b)\big) &&
\end{flalign*}

\subsubsection*{ب}
خیر. این نتیجه تنها می‌گوید ضرب داخلیِ دو بردار کدگذاری مکانیِ خام فقط به اختلاف $a-b$ بستگی دارد؛ یعنی خودِ کدگذاری ساختاری متناسب با فاصلهٔ نسبی دارد. اما این به معنای استفادهٔ صرفاً نسبی از اطلاعات مکانی در کل شبکه نیست:

نخست آنکه کدگذاری مکانی به‌صورت جمع با \lr{embedding} توکن به ورودی افزوده می‌شود، پس هر موقعیت یک کدِ مطلق دارد و اطلاعات مطلق در دسترس شبکه قرار می‌گیرد. دوم آنکه امتیاز توجهِ واقعی برابر $p_a^T p_b$ نیست؛ بلکه پرس‌وجو و کلید نگاشتِ خطیِ مجموع (\lr{embedding} + \lr{PE}) هستند و با محتوا نیز برهم‌کنش دارند، و خاصیتِ نسبی‌بودنِ $p_a^T p_b$ پس از اعمال $W_Q, W_K$ و افزودنِ محتوا لزوماً پابرجا نمی‌ماند. سوم آنکه لایه‌های بعدی می‌توانند از الگوی سینوسی موقعیتِ مطلق را بازسازی و استفاده کنند.

پس نتیجهٔ بند قبل صرفاً یک ویژگیِ ضرب‌داخلیِ کدهای مکانیِ خالص است، نه اثباتِ نسبی‌بودنِ استفاده از مکان در سراسر مدل.

\section*{پرسش ۴}

\subsection*{الف}
داریم $d_{\text{model}}=768$، اندازهٔ واژگان $V=10000$ و بیشینهٔ طول دنباله $L=1000$. در شمارش، بایاس‌ها و دو \lr{LayerNorm} هر بلوک را نیز لحاظ می‌کنیم.

embedding:
\begin{flalign*}
\underbrace{10000\times 768}_{\text{توکن}} + \underbrace{1000\times 768}_{\text{مکانی}} = 7{,}680{,}000 + 768{,}000 = 8{,}448{,}000 &&
\end{flalign*}

یک بلوک دیکدر شامل \lr{self-attention} نقاب‌دار و \lr{MLP} است. چون $h\cdot d_k = 4\times 192 = 768 = d_{\text{model}}$، هر یک از نگاشت‌های $Q,K,V,O$ ماتریسی $768\times768$ است:
\begin{flalign*}
\text{attention} &= 4\times(768\times768 + 768) = 4\times 590{,}592 = 2{,}362{,}368 \\
\text{MLP} &= (768\times3000 + 3000) + (3000\times768 + 768) = 2{,}307{,}000 + 2{,}304{,}768 = 4{,}611{,}768 \\
\text{دو LayerNorm} &= 2\times(2\times768) = 3{,}072 \\
\text{یک بلوک} &= 2{,}362{,}368 + 4{,}611{,}768 + 3{,}072 = 6{,}977{,}208 &&
\end{flalign*}
(در صورت صرف‌نظر از بایاس و \lr{LayerNorm}: $4\times768^2 + 2\times768\times3000 = 6{,}967{,}296$.)

لایهٔ خروجیِ پیش از \lr{softmax} که $d_{\text{model}}\to V$ را نگاشت می‌کند:
\begin{flalign*}
768\times 10000 + 10000 = 7{,}690{,}000 &&
\end{flalign*}
(بدون بایاس $7{,}680{,}000$؛ در صورت اشتراک وزن با \lr{embedding} ورودی، پارامتر تازه‌ای افزوده نمی‌شود.)

\subsection*{ب}
یک لایهٔ RNN ساده با $h_t = \tanh(W_{xh}x_t + W_{hh}h_{t-1} + b)$ و ورودی و حالت پنهانِ هر دو با بعد ۷۶۸:
\begin{flalign*}
\underbrace{768\times768}_{W_{xh}} + \underbrace{768\times768}_{W_{hh}} + \underbrace{768}_{b} = 1{,}179{,}648 + 768 = 1{,}180{,}416 &&
\end{flalign*}
که بسیار کمتر از پارامترهای یک بلوک ترنسفرمر است.

\subsection*{ج}
LSTM چهار دروازه (ورودی، فراموشی، خروجی و کاندیدای سلول) دارد و هر دروازه ساختاری همانند RNN بالا:
\begin{flalign*}
4\times\big(768\times768 + 768\times768 + 768\big) = 4\times 1{,}180{,}416 = 4{,}721{,}664 &&
\end{flalign*}

\subsection*{د}
کانولوشن یک‌بعدی با ۷۶۸ کانال ورودی، ۷۶۸ فیلتر و اندازهٔ کرنل ۳:
\begin{flalign*}
\underbrace{3\times768\times768}_{\text{وزن}} + \underbrace{768}_{\text{بایاس}} = 1{,}769{,}472 + 768 = 1{,}770{,}240 &&
\end{flalign*}
ایراد اساسی این است که میدان دیدِ این لایه محلی و محدود به پنجره‌ای ۳-تایی است؛ یعنی یک لایهٔ کانولوشن تنها وابستگیِ میانِ توکن‌های بسیار نزدیک را می‌بیند و قادر به مدل‌کردنِ مستقیمِ وابستگی‌های دوربرد در دنباله نیست. برخلاف توجه (با میدان دیدِ سراسری) و RNN/LSTM (با انتشار حالت در کل دنباله)، برای پوشش وابستگی‌های بلندمدت باید شمار زیادی لایهٔ کانولوشن روی هم انباشته شوند.